\documentclass[a4paper,12pt]{article}
\renewcommand{\rmdefault}{nmr}
\renewcommand{\sfdefault}{nmss}
\usepackage[T2A]{fontenc}
\usepackage[utf8]{inputenc}
\usepackage[russian]{babel}
\usepackage{amsmath,amssymb,mathtools}
\usepackage{nm}
\renewcommand{\familydefault}{\sfdefault}
\usepackage{icomma}
\usepackage[a4paper,top=19mm,bottom=19mm,left=20mm,right=20mm,headheight=15pt]{geometry}
\usepackage{microtype}
\usepackage{tikz}
\usetikzlibrary{calc,arrows.meta,positioning,shapes.geometric}
\usepackage{tabularx,booktabs}
\usepackage{enumitem}
\usepackage{parskip}
\usepackage{setspace}
\usepackage{needspace}
\usepackage{xcolor}
\usepackage{fancyhdr}
\usepackage{titlesec}
\usepackage{hyperref}
\definecolor{accent}{HTML}{007F80}
\definecolor{darkaccent}{HTML}{004D4D}
\definecolor{textgray}{HTML}{333333}
\definecolor{softgray}{HTML}{7C8689}
\hypersetup{hidelinks,pdftitle={Чек-лист. Финансовая математика. 20.10.2026},pdfauthor={Артём Александрович Лёвин}}
\onehalfspacing
\setlength{\parindent}{0pt}
\setlength{\parskip}{5pt}
\titleformat{\section}{\large\bfseries\color{darkaccent}}{\thesection.}{0.6em}{}
\titleformat{\subsection}{\normalsize\bfseries\color{darkaccent}}{\thesubsection.}{0.6em}{}
\titlespacing*{\section}{0pt}{15pt}{5pt}
\titlespacing*{\subsection}{0pt}{11pt}{3pt}
\setlist[itemize]{leftmargin=*,topsep=3pt,itemsep=1pt}
\setlist[enumerate]{leftmargin=*,topsep=3pt,itemsep=5pt}
\newcommand{\checkitem}{\(\square\)\enspace}
\newcommand{\dd}{\displaystyle}
\pagestyle{fancy}
\fancyhf{}
\fancyhead[L]{\small\color{softgray}Финансовая математика}
\fancyhead[R]{\small\color{softgray}20.10.2026}
\fancyfoot[C]{\small\thepage}
\renewcommand{\headrulewidth}{0.3pt}
\renewcommand{\footrulewidth}{0pt}
\begin{document}

% Титульная страница
\thispagestyle{empty}
\vspace*{42mm}
\begin{center}
{\color{accent}\large МАТЕМАТИКА \quad / \quad ЕГЭ, профильный уровень}\par
\vspace{18mm}
{\Huge\bfseries Чек-лист}\par
\vspace{8mm}
{\LARGE\bfseries\color{darkaccent}Финансовая математика}\par
\vspace{4mm}
{\Large Вклады, кредиты и переплата}\par
\vspace{17mm}
{\color{accent}\rule{68mm}{0.8pt}}\par
\vspace{16mm}
{\large 20 октября 2026 года}\par
\vspace{13mm}
{\normalsize\color{textgray}Лёвин Артём Александрович}\\[5pt]
{\normalsize\bfseries эксклюзивно для Тимофея Васильченко}
\end{center}
\vfill
\noindent\begin{tikzpicture}
\draw[accent!55,line width=1.0pt] (0,0) .. controls (4.9,0.9) and (11.7,-0.4) .. (16.5,0.2);
\end{tikzpicture}
\vspace*{8mm}
\clearpage

\section*{Что нужно уметь после занятия}
\checkitem Отличать начисленные проценты от суммы с процентами.\\
\checkitem Рассчитывать вклад с капитализацией и пополнением.\\
\checkitem Узнавать аннуитетную и дифференцированную схемы по условию.\\
\checkitem Строить таблицу долга для смешанной или произвольной схемы.\\
\checkitem Вычислять сумму платежей и процент переплаты.\\
\checkitem Сравнивать два способа погашения при одинаковых условиях.

\section{Один процентный коэффициент --- две величины}
Пусть \(r\) --- годовая процентная ставка, \(p=r/100\) --- ставка в долях единицы, а
\[
\boxed{k=1+p=1+\frac r{100}}.
\]

Если на счёте или в долге находится сумма \(S\), то:
\[
\underbrace{Sp}_{\text{только начисленные проценты}},\qquad
\underbrace{Sk}_{\text{сумма вместе с процентами}}.
\]
\textbf{Пример.} При ставке \(20\%\) имеем \(p=0,2\), \(k=1,2\).
На \(100\) тыс. рублей начислят \(20\) тыс., а сумма после начисления
составит \(120\) тыс. рублей.

\textbf{Не перепутайте:} в колонке «начисленные проценты» стоит \(Sp\),
в колонке «сумма с процентами» --- \(Sk\).

\section{Вклады с ежегодной капитализацией}
Если проценты в конце года остаются на вкладе, то через \(n\) лет без пополнений
\[
\boxed{S_n=Sk^n}.
\]
При дополнительном внесении суммы \(a\) после первого начисления процентов
(и ещё одном начислении в конце второго года):
\[
\boxed{S_2=(Sk+a)k=Sk^2+ak}.
\]
При нескольких пополнениях \(a_1,a_2,\ldots,a_{n-1}\), внесённых после очередного годового начисления,
\[
S_n=Sk^n+a_1k^{n-1}+a_2k^{n-2}+\cdots+a_{n-1}k.
\]
При снятии средств соответствующий член берётся со знаком «минус».

\Needspace{15\baselineskip}
\subsection*{Пример из занятия}
В январе 2028 года внесли \(950\) тыс. рублей, в январе 2029 года добавили
\(550\) тыс. рублей. Проценты начисляются в конце года. В январе 2030 года
на счёте \(1807,68\) тыс. рублей. Найти годовую ставку.
\[
950k^2+550k=1807,68.
\]
Подстановка \(k=1,12\) даёт
\[
950\cdot1,12^2+550\cdot1,12=1191,68+616=1807,68.
\]
Поскольку \(k\ge1\) и левая часть возрастает по \(k\), найденное
значение единственно: \(r=(1,12-1)\cdot100\%=\boxed{12\%}\).

\section{Универсальная таблица кредита}
Обозначим \(B_i\) --- долг \emph{перед} очередным начислением процентов,
\(X_i\) --- платёж после начисления, \(B_{i+1}\) --- остаток долга
после платежа. Тогда
\[
\boxed{\text{Долг с процентами}=kB_i,\qquad
X_i=kB_i-B_{i+1}.}
\]
В последнем периоде кредит погашается полностью: \(B_{n+1}=0\).

\renewcommand{\arraystretch}{1.3}
\begin{center}
\begin{tabularx}{\linewidth}{@{}lXXX@{}}
\toprule
Этап & Долг до процентов & Долг с процентами & Платёж / остаток \\
\midrule
\(i\)-й год & \(B_i\) & \(kB_i\) & \(X_i=kB_i-B_{i+1}\) \\
Последний год & \(B_n\) & \(kB_n\) & \(X_n=kB_n\) \\
\bottomrule
\end{tabularx}
\end{center}

\textbf{Правило чтения условия.} Сначала установите, что именно задано:
равный платёж, равная часть \emph{основного долга} или график \emph{остатка долга}.
Это разные требования.

\section{Дифференцированная схема}
\textbf{Признак:} «долг уменьшается на одну и ту же сумму», «основной долг
гасится равными долями».

При кредите \(S\) на \(n\) лет основная часть каждого платежа равна \(S/n\).
Поэтому для \(i=1,\ldots,n\):
\[
\boxed{X_i=\frac Sn+p\left(S-\frac{(i-1)S}{n}\right).}
\]
Сумма всех платежей:
\[
\boxed{T_{\text{диф}}=S+Sp\frac{n+1}{2}.}
\]
Это следует из того, что проценты начисляются на остатки
\(S,\frac{n-1}{n}S,\ldots,\frac1nS\).

\textbf{Пример.} Кредит \(1\) млн рублей на \(2\) года, общая сумма выплат
\(1255\) тыс. рублей. Из формулы
\[
1255=1000+1000p\frac{3}{2}
\quad\Longrightarrow\quad p=0,17
\quad\Longrightarrow\quad\boxed{r=17\%}.
\]

\section{Аннуитетная схема}
\textbf{Признак:} «погашение \emph{равными платежами}».
Если каждый год после начисления процентов вносится одна сумма \(x\), то
\[
Sk^n-xk^{n-1}-xk^{n-2}-\cdots-x=0.
\]
Следовательно,
\[
\boxed{x=\frac{Sk^n}{k^{n-1}+k^{n-2}+\cdots+k+1}},
\qquad \boxed{T_{\text{анн}}=nx}.
\]
Формула работает и при нулевой ставке: тогда \(k=1\), \(x=S/n\).

\textbf{Пример.} Кредит \(496,5\) тыс. рублей на \(3\) года под \(10\%\):
\[
x=\frac{496,5\cdot1,1^3}{1,1^2+1,1+1}
=\frac{660,8415}{3,31}=199,65\ \text{тыс. руб.}
\]
\[
\boxed{T_{\text{анн}}=3\cdot199,65=598,95\ \text{тыс. руб.}}
\]
\clearpage

\section{Смешанная схема: остаток уменьшается на процент}
\textbf{Признак:} «долг к следующему июлю на \(20\%\) меньше долга к предыдущему июлю».
Это условие на \emph{остаток}, а не на равенство платежей или равенство
погашаемых частей основного долга. Удобнее построить таблицу.

\textbf{Пример.} В июле 2029 года взяли \(1,4\) млн рублей на \(4\) года.
В январе долг увеличивается на \(20\%\), затем совершается платёж.
Остатки в июле 2030--2032 годов каждый раз составляют \(80\%\) предыдущего.
К июлю 2033 года долг полностью погашен.

\smallskip
\begin{center}
\small
\begin{tabular}{@{}ccccc@{}}
\toprule
Год & До начисления & После начисления & Платёж & Остаток \\
\midrule
2030 & 1,4 & 1,68 & 0,56 & 1,12 \\
2031 & 1,12 & 1,344 & 0,448 & 0,896 \\
2032 & 0,896 & 1,0752 & 0,3584 & 0,7168 \\
2033 & 0,7168 & 0,86016 & 0,86016 & 0 \\
\bottomrule
\end{tabular}
\end{center}
\begin{center}\footnotesize Все суммы в таблице --- в миллионах рублей.\end{center}

Искомый последний платёж:
\[
X_{2033}=1,4\cdot0,8^3\cdot1,2=0,86016\ \text{млн руб.}
=\boxed{860\,160\ \text{руб.}}
\]
\textbf{Внимание к единицам!} Если расчёты выполнены в миллионах рублей,
для ответа в рублях результат надо умножить на \(1\,000\,000\).

\section{Произвольный график и переплата}
По условию полугодового кредита 15 января выдали сумму \(S\).
Остатки долга к 15-му числу последовательных месяцев:
\(100\%,90\%,80\%,70\%,60\%,50\%,0\%\) от \(S\).
В конце каждого месяца перед очередным платежом остаток увеличивается на \(5\%\).

Возьмём \(S=100\) условных единиц. Каждый платёж рассчитывается как
«долг после процентов минус указанный следующий остаток».
\begin{center}
\small
\begin{tabular}{@{}ccccc@{}}
\toprule
Месяц & Начало & После \(5\%\) & След. остаток & Платёж \\
\midrule
1 & 100 & 105 & 90 & 15 \\
2 & 90 & 94,5 & 80 & 14,5 \\
3 & 80 & 84 & 70 & 14 \\
4 & 70 & 73,5 & 60 & 13,5 \\
5 & 60 & 63 & 50 & 13 \\
6 & 50 & 52,5 & 0 & 52,5 \\
\bottomrule
\end{tabular}
\end{center}
Всего выплачено \(T=122,5\) условной единицы, то есть \(122,5\%\) суммы кредита.

\textbf{Абсолютная переплата:} \(P=T-S\).\quad
\textbf{Процент переплаты:}
\[
\boxed{q=\frac{T-S}{S}\cdot100\%.}
\]
В этом примере \(q=(122,5-100)/100\cdot100\%=\boxed{22,5\%}\).

\section{Как сравнивать две схемы}
Сравнивать выгодность следует по \textbf{общей сумме платежей}\,
(если размер кредита, ставка и сроки совпадают).

\textbf{Пример.} \(S=100\) тыс. рублей, \(r=20\%\), \(n=3\) года,
платежи ежегодные.
\[
T_{\text{диф}}=100+100\cdot0,2\cdot\frac{4}{2}=140\ \text{тыс. руб.},
\]
\[
T_{\text{анн}}=3\cdot\frac{100\cdot1,2^3}{1+1,2+1,2^2}
\approx142,42\ \text{тыс. руб.}
\]
При этих условиях дифференцированная схема дешевле примерно на
\(\boxed{2,42\ \text{тыс. руб.}}\)

\textbf{Самопроверка:} сумма всех выплат должна быть не меньше первоначального
кредита при неотрицательной ставке; остаток после последнего платежа --- нуль.

\clearpage
\thispagestyle{fancy}
\section*{Алгоритм. Как выбрать способ решения задачи}
\vspace{10mm}
\begin{center}
\begin{tikzpicture}[
  >=Stealth,
  line/.style={-{Stealth[length=3mm,width=2mm]},thick,draw=darkaccent},
  start/.style={ellipse,draw=darkaccent,line width=0.9pt,align=center,minimum width=58mm,minimum height=12mm,fill=accent!5},
  dec/.style={diamond,aspect=2.25,draw=darkaccent,line width=0.9pt,align=center,text width=47mm,inner sep=1pt,minimum height=22mm},
  final/.style={ellipse,draw=accent,line width=0.9pt,align=center,text width=52mm,minimum height=14mm},
  tag/.style={font=\small,fill=white,inner sep=1.5pt}]
\node[start] (begin) at (0,0) {Прочитайте условие кредита};
\node[dec] (equal) at (0,-3.1) {Все платежи равны?};
\node[dec] (diff) at (0,-6.8) {Основной долг гасится равными долями?};
\node[dec] (balance) at (0,-10.5) {Заданы остатки долга по датам?};
\node[final] (a) at (7.2,-3.1) {Аннуитет: формула или таблица};
\node[final] (d) at (7.2,-6.8) {Дифференцированная схема};
\node[final] (b) at (7.2,-10.5) {Смешанная или произвольная схема: таблица};
\node[final] (t) at (0,-13.8) {Универсальная таблица: \(X_i=kB_i-B_{i+1}\)};
\draw[line] (begin) -- (equal);
\draw[line] (equal.east) -- node[tag,above]{да} (a.west);
\draw[line] (equal.south) -- node[tag,right]{нет} (diff.north);
\draw[line] (diff.east) -- node[tag,above]{да} (d.west);
\draw[line] (diff.south) -- node[tag,right]{нет} (balance.north);
\draw[line] (balance.east) -- node[tag,above]{да} (b.west);
\draw[line] (balance.south) -- node[tag,right]{нет} (t.north);
\end{tikzpicture}
\end{center}
\clearpage

\section*{Самостоятельная тренировка \textemdash{} Путь А}
\textbf{10 заданий, от базовых к более сложным.}
При расчётах указывайте единицы измерения и различайте остаток долга,
начисленные проценты и платёж.
\begin{enumerate}[label=\textbf{\arabic*.},itemsep=9pt]
\item Ставка вклада равна \(25\%\) годовых. Найдите \(p\), \(k\) и сумму на счёте через год, если первоначально внесли \(80\) тыс. рублей.
\item На вклад внесли \(200\) тыс. рублей под \(10\%\) годовых. Проценты капитализируются, пополнений нет. Сколько будет на счёте через \(2\) года?
\item На вклад внесли \(300\) тыс. рублей под \(10\%\) годовых. После первого начисления процентов внесли ещё \(50\) тыс. рублей. Найдите сумму после второго начисления.
\item На вклад положили \(500\) тыс. рублей на \(2\) года. После первого года добавили \(200\) тыс. рублей. После второго начисления на счёте оказалось \(825\) тыс. рублей. Определите годовую ставку.
\item Дифференцированный кредит \(300\) тыс. рублей выдан на \(3\) года под \(10\%\) годовых. Найдите платёж второго года и общую сумму платежей.
\item Дифференцированный кредит \(400\) тыс. рублей рассчитан на \(3\) года. Общая сумма выплат \(480\) тыс. рублей. Найдите годовую ставку.
\item Аннуитетный кредит \(210\) тыс. рублей выдан на \(2\) года под \(10\%\) годовых. Найдите каждый равный ежегодный платёж и общую сумму выплат.
\item Кредит \(500\) тыс. рублей выдан на \(2\) года под \(20\%\) годовых. После первого начисления процентов и платежа остаток составляет \(60\%\) первоначального долга. После второго начисления кредит погашают. Найдите оба платежа и переплату в рублях.
\item Кредит погашают за \(3\) месяца. Остатки долга перед начислением процентов: \(100\%\), затем \(80\%\), затем \(30\%\), после последнего платежа \(0\%\). Каждый месяц перед платежом долг увеличивается на \(10\%\). Найдите все платежи в процентах от первоначального кредита и процент переплаты.
\item Кредит \(331\) тыс. рублей выдан на \(3\) года под \(10\%\) годовых. Сравните общую сумму выплат по аннуитетной и дифференцированной схемам. Какая дешевле и на сколько рублей?
\end{enumerate}
\clearpage
\section*{Ответы для самопроверки}
\renewcommand{\arraystretch}{1.52}
\begin{tabularx}{\linewidth}{@{}cX@{}}
\toprule
№ & Ответ \\
\midrule
1 & \(p=0,25\); \(k=1,25\); \(100\) тыс. руб. \\
2 & \(242\) тыс. руб. \\
3 & \(418\) тыс. руб. \\
4 & \(10\%\). \\
5 & Второй платёж --- \(120\) тыс. руб.; всего --- \(360\) тыс. руб. \\
6 & \(10\%\). \\
7 & Ежегодный платёж --- \(121\) тыс. руб.; всего --- \(242\) тыс. руб. \\
8 & \(300\) тыс. руб. и \(360\) тыс. руб.; переплата \(160\,000\) руб. \\
9 & \(30\%\), \(58\%\), \(33\%\) от суммы кредита; переплата \(21\%\). \\
10 & Аннуитет: \(399\,300\) руб.; дифференцированная схема: \(397\,200\) руб.; дифференцированная выгоднее на \(2\,100\) руб. \\
\bottomrule
\end{tabularx}

\vspace{10mm}
\textbf{Проверочные формулы:}
\[
\begin{gathered}
 k=1+\frac r{100},\qquad S_2=Sk^2+ak,\\[4pt]
 X_i=kB_i-B_{i+1},\qquad
 T_{\text{диф}}=S+Sp\frac{n+1}{2},\\[4pt]
 x=\frac{Sk^n}{\sum_{j=0}^{n-1}k^j},\qquad
 q=\frac{T-S}{S}\cdot100\%.
\end{gathered}
\]
\end{document}